%beamer

% Start with video if you like it and can get the sound working well:
% https://www.youtube.com/watch?v=Dy0hJWltsyE

% Comment/uncomment this line to toggle handout mode
\newcommand{\handout}{}

%% Beamer-Klasse im korrekten Modus
\ifdefined \handout
\documentclass[handout]{beamer} % Handout mode
\else
\documentclass{beamer}
\fi

%% UTF-8-Encoding
\usepackage[utf8]{inputenc}

% % \bigtimes abgeschrieben von http://tex.stackexchange.com/questions/14386/importing-a-single-symbol-from-a-different-font
% \DeclareFontFamily{U}{mathx}{\hyphenchar\font45}
% \DeclareFontShape{U}{mathx}{m}{n}{
%       <5> <6> <7> <8> <9> <10> gen * mathx
%       <10.95> mathx10 <12> <14.4> <17.28> <20.74> <24.88> mathx12
%       }{}
% \DeclareSymbolFont{mathx}{U}{mathx}{m}{n}
% \DeclareMathSymbol{\bigtimes}{\mathop}{mathx}{161}

\RequirePackage{xcolor}

\def\9{\square}
%\def\9{\blank}

% f"ur Aussagenlogik
\colorlet{alcolor}{blue}
\RequirePackage{tikz}
\usetikzlibrary{arrows.meta}
\newcommand{\alimpl}{\mathrel{\tikz[x={(0.1ex,0ex)},y={(0ex,0.1ex)},>={Classical TikZ Rightarrow[]}]{\draw[alcolor,->,line width=0.7pt,line cap=round] (0,0) -- (15,0);\path (0,-6);}}}
\newcommand{\aleqv}{\mathrel{\tikz[x={(0.1ex,0ex)},y={(0ex,0.1ex)},>={Classical TikZ Rightarrow[]}]{\draw[alcolor,<->,line width=0.7pt,line cap=round] (0,0) -- (18,0);\path (0,-6);}}}
\newcommand{\aland}{\mathbin{\raisebox{-0.6pt}{\rotatebox{90}{\texttt{\color{alcolor}\char62}}}}}
\newcommand{\alor}{\mathbin{\raisebox{-0.8pt}{\rotatebox{90}{\texttt{\color{alcolor}\char60}}}}}
%\newcommand{\ali}[1]{_{\mathtt{\color{alcolor}#1}}}
\newcommand{\alv}[1]{\mathtt{\color{alcolor}#1}}
\newcommand{\alnot}{\mathop{\tikz[x={(0.1ex,0ex)},y={(0ex,0.1ex)}]{\draw[alcolor,line width=0.7pt,line cap=round,line join=round] (0,0) -- (10,0) -- (10,-4);\path (0,-8) ;}}}
\newcommand{\alP}{\alv{P}} %ali{#1}}
%\newcommand{\alka}{\negthinspace\hbox{\texttt{\color{alcolor}(}}}
\newcommand{\alka}{\negthinspace\text{\texttt{\color{alcolor}(}}}
%\newcommand{\alkz}{\texttt{\color{alcolor})}}\negthinspace}
\newcommand{\alkz}{\text{\texttt{\color{alcolor})}}\negthinspace}
\newcommand{\AAL}{A_{AL}}
\newcommand{\LAL}{\hbox{\textit{For}}_{AL}}
\newcommand{\AxAL}{\hbox{\textit{Ax}}_{AL}}
\newcommand{\AxEq}{\hbox{\textit{Ax}}_{Eq}}
\newcommand{\AxPL}{\hbox{\textit{Ax}}_{PL}}
\newcommand{\AALV}{\hbox{\textit{Var}}_{AL}}
\newcommand{\MP}{\hbox{\textit{MP}}}
\newcommand{\GEN}{\hbox{\textit{GEN}}}
\newcommand{\W}{\ensuremath{\hbox{\textbf{w}}}\xspace}
\newcommand{\F}{\ensuremath{\hbox{\textbf{f}}}\xspace}
\newcommand{\WF}{\ensuremath{\{\W,\F\}}\xspace}
\newcommand{\val}{\hbox{\textit{val}}}
\newcommand{\valDIb}{\val_{D,I,\beta}}

\newcommand*{\from}{\colon}

% die nachfolgenden Sachen angepasst an cmtt
\newlength{\ttquantwd}
\setlength{\ttquantwd}{1ex}
\newlength{\ttquantht}
\setlength{\ttquantht}{6.75pt}
\def\plall{%
  \tikz[line width=0.67pt,line cap=round,line join=round,baseline=(B),alcolor] {
    \draw (-0.5\ttquantwd,\ttquantht) -- node[coordinate,pos=0.4] (lll){} (-0.25pt,-0.0pt) -- (0.25pt,-0.0pt) -- node[coordinate,pos=0.6] (rrr){} (0.5\ttquantwd,\ttquantht);
    \draw (lll) -- (rrr);
    \coordinate (B) at (0,-0.35pt);
  }%
}
\def\plexist{%
  \tikz[line width=0.67pt,line cap=round,line join=round,baseline=(B),alcolor] {
    \draw (-0.9\ttquantwd,\ttquantht) -- (0,\ttquantht) -- node[coordinate,pos=0.5] (mmm){} (0,0) --  (-0.9\ttquantwd,0);
    \draw (mmm) -- ++(-0.75\ttquantwd,0);
    \coordinate (B) at (0,-0.35pt);
  }\ensuremath{\,}%
}
\let\plexists=\plexist
\newcommand{\NT}[1]{\ensuremath{\langle\mathrm{#1} \rangle}}

\newcommand{\CPL}{\text{\itshape Const}_{PL}}
\newcommand{\FPL}{\text{\itshape Fun}_{PL}}
\newcommand{\RPL}{\text{\itshape Rel}_{PL}}
\newcommand{\VPL}{\text{\itshape Var}_{PL}}
\newcommand{\ATer}{A_{\text{\itshape Ter}}}
\newcommand{\ARel}{A_{\text{\itshape Rel}}}
\newcommand{\AFor}{A_{\text{\itshape For}}}
\newcommand{\LTer}{L_{\text{\itshape Ter}}}
\newcommand{\LRel}{L_{\text{\itshape Rel}}}
\newcommand{\LFor}{L_{\text{\itshape For}}}
\newcommand{\NTer}{N_{\text{\itshape Ter}}}
\newcommand{\NRel}{N_{\text{\itshape Rel}}}
\newcommand{\NFor}{N_{\text{\itshape For}}}
\newcommand{\PTer}{P_{\text{\itshape Ter}}}
\newcommand{\PRel}{P_{\text{\itshape Rel}}}
\newcommand{\PFor}{P_{\text{\itshape For}}}

\newcommand{\plka}{\alka}
\newcommand{\plkz}{\alkz}
%\newcommand{\plka}{\plfoo{(}}
%\newcommand{\plkz}{\plfoo{)}}
\newcommand{\plcomma}{\hbox{\texttt{\color{alcolor},}}}
\newcommand{\pleq}{{\color{alcolor}\,\dot=\,}}

% MODIFIED (DJ)
% previously: \newcommand{\plfoo}[1]{\mathtt{\color{alcolor}#1}}
\newcommand{\plfoo}[1]{\texttt{\color{alcolor}#1}}

\newcommand{\plc}{\plfoo{c}}
\newcommand{\pld}{\plfoo{d}}
\newcommand{\plf}{\plfoo{f}}
\newcommand{\plg}{\plfoo{g}}
\newcommand{\plh}{\plfoo{h}}
\newcommand{\plx}{\plfoo{x}}
\newcommand{\ply}{\plfoo{y}}
\newcommand{\plz}{\plfoo{z}}
\newcommand{\plR}{\plfoo{R}}
\newcommand{\plS}{\plfoo{S}}

\newcommand{\bv}{\mathrm{bv}}
\newcommand{\fv}{\mathrm{fv}}

%\newcommand{\AxAL}{\hbox{\textit{Ax}}_{AL}}
%\newcommand{\AALV}{\hbox{\textit{Var}}_{AL}}

%\renewcommand{\#}[1]{\literal{#1}}
\newcommand{\A}{\mathcal{A}}
\newcommand{\Adr}{\text{Adr}}
\newcommand{\ar}{\mathrm{ar}}
\newcommand{\ascii}[1]{\literal{\char#1}}
%\newcommand{\assert}[1]{\text{/\!\!/\ } #1}
\newcommand{\assert}[1]{\colorbox{black!7!white}{\ensuremath{\{\;#1\;\}}}}
\newcommand{\Assert}[1]{$\langle$\textit{#1}$\rangle$}
\newcommand{\B}{\mathcal{B}}
\newcommand{\bfmod}{\mathbin{\kw{ mod }}}
\newcommand{\bb}{{\text{bb}}}
\def\bottom{\hbox{\small$\pmb{\bot}$}}
\newcommand{\card}[1]{|#1|}
%\newcommand{\cod}{\mathop{\text{cod}}}  % ist in thwmathabbrevs
\newcommand{\Conf}{\mathcal{C}}
\newcommand{\define}[1]{\emph{#1}}
%\renewcommand{\dh}{d.\,h.\@\xspace}
%\newcommand{\Dh}{D.\,h.\@\xspace}
%\newcommand{\engl}[1]{engl.\xspace\emph{#1}}
\newcommand{\eps}{\varepsilon}
%\newcommand{\evtl}{evtl.\@\xspace}
\newcommand{\fbin}{\text{bin}}
\newcommand{\finv}{\text{inv}}
\newcommand{\fnum}{\text{num}}
\newcommand{\fNum}{{\text{Num}}}
\newcommand{\frepr}{\text{repr}}
\newcommand{\fRepr}{\text{Repr}}
\newcommand{\fZkpl}{\text{Zkpl}}
\newcommand{\fLen}{\text{Len}}
\newcommand{\fsem}{\text{sem}}
\providecommand{\fspace}{\mathord{\text{space}}}
\providecommand{\fSpace}{\mathord{\text{Space}}}
\providecommand{\ftime}{\mathord{\text{time}}}
\providecommand{\fTime}{\mathord{\text{Time}}}
\newcommand{\fTrans}{\text{Trans}}
\newcommand{\fVal}{\text{Val}}

% Modified (DJ)
\newcommand{\Val}{\text{Val}}

%\def\G{\mathbb{Z}}
\newcommand{\HT}[1]{\normalfont\textsc{HT-#1}}
\newcommand{\htr}[3]{\{#1\}\;#2\; \{#3\}}
\newcommand{\Id}{\text{I}}
%\newcommand{\ie}{i.\,e.\@\xspace}
\newcommand{\instr}[2]{\texttt{#1}\ \textit{#2}}
\newcommand{\Instr}[2]{\texttt{#1}\ \textrm{#2}}
\newcommand{\instrr}[3]{\texttt{#1}\ \textit{#2}\texttt{(#3)}}
\newcommand{\Instrr}[3]{\texttt{#1}\ \textrm{#2}\texttt{(#3)}}
\newcommand{\io}{\!\mid\!}
\usepackage{KITcolors}
\newcommand{\literal}[1]{\hbox{\textcolor{blue!95!white}{\textup{\texttt{\scalebox{1.11}{#1}}}}}}
%\newcommand{\literal}[1]{\hbox{\textcolor{KITblue!80!black}{\textup{\texttt{#1}}}}}
\def\kasten#1{\leavevmode\literal{\setlength{\fboxsep}{1pt}\fbox{\vrule  width 0pt height 1.5ex depth 0.5ex #1}}}
\newcommand{\kw}[1]{\ensuremath{\mathbf{#1}}}
\newcommand{\lang}[1]{\ensuremath{\langle#1\rangle}}
%\newcommand{\maw}{m.\,a.\,w.\@\xspace}
%\newcommand{\MaW}{M.\,a.\,w.\@\xspace}
\newcommand{\mdefine}[2][FOOBAR]{\define{#2}\def\foobar{FOOBAR}\def\optarg{#1}\ifx\foobar\optarg\def\optarg{#2}\fi\graffito{\optarg}}
\newcommand{\meins}{\rotatebox[origin=c]{180}{1}}
\newcommand{\Mem}{\text{Mem}}
\newcommand{\memread}{\text{memread}}
\newcommand{\memwrite}{\text{memwrite}}
\providecommand{\meta}[1]{\ensuremath{\langle}\textit{#1}\ensuremath{\rangle}}
%\newcommand{\N}{\mathbb{N}}
\newcommand{\NP}{\mathbf{NP}}
\newcommand{\Nadd}{N_{\text{add}}}
\newcommand{\Nmult}{N_{\text{mult}}}
\newcommand{\Oh}[1]{O\left(#1\right)}
\newcommand{\Om}[1]{\Omega\left(#1\right)}
\newcommand{\personname}[1]{\textsc{#1}}
\newcommand{\regname}[1]{\texttt{#1}}
\newcommand{\mima}{\textsc{Mima}\xspace}
\newcommand{\mimax}{\textsc{Mima-X}\xspace}

\def\Pclass{\text{\bfseries P}}
\def\PSPACE{\text{\bfseries PSPACE}}

\newcommand{\SPush}{\text{push}}
\newcommand{\SPop}{\text{pop}}
\newcommand{\SPeek}{\text{peek}}
\newcommand{\STop}{\text{top}}
\newcommand{\STos}{\text{\itshape tos}}
\newcommand{\SBos}{\text{\itshape bos}}

%\newcommand{\R}{\mathbb{R}}
\newcommand{\Rnullplus}{\R_0^{+}}
\newcommand{\Rplus}{\R_{+}}
\newcommand{\resp}{resp.\@\xspace}
\newcommand{\Sem}{\text{Sem}}
\newcommand{\sgn}{\mathop{\text{sgn}}}
\newcommand{\sqbox}{\mathop{\raisebox{-6.2pt}{\hbox{\hbox to 0pt{$^{^{\sqcap}}$\hss}$^{^{\sqcup}}$}}}}
\newcommand{\sqleq}{\sqsubseteq}
\newcommand{\sqgeq}{\sqsupseteq}
\newcommand{\Th}[1]{\Theta\left(#1\right)}
%\newcommand{\usw}{usw.\@\xspace}
\newcommand{\V}[1]{\hbox{\textit{#1}}}
\newcommand{\x}{\times}
\newcommand{\ZK}{\mathbb{K}}
%\newcommand{\Z}{\mathbb{Z}}
%\newcommand{\zB}{z.\,B.\@\xspace}
%\newcommand{\ZB}{Z.\,B.\@\xspace}
% \newcommand{\bb}{{\text{bb}}}
% \def\##1{\hbox{\textcolor{darkblue}{\texttt{#1}}}}
% \def\A{\mathcal{A}}
% \newcommand{\0}{\#0}
% \newcommand{\1}{\#1}
% \newcommand{\Obj}{\text{Obj}}
% \newcommand{\start}{\mathop{\text{start}}}
% \newcommand{\compactlist}{\addtolength{\itemsep}{-\parskip}}
% \newcommand{\fval}{\text{val}}
% \newcommand{\lang}[1]{\ensuremath{\langle#1\rangle}}
% \newcommand{\io}{\!\mid\!}
% \def\sqbox{\mathop{\raisebox{-6.2pt}{\hbox{\hbox to 0pt{$^{^{\sqcap}}$\hss}$^{^{\sqcup}}$}}}}
% \def\sqleq{\sqsubseteq}
% \def\sqgeq{\sqsupseteq}
\def\Td{T_{\overline{d}}}
% \newcommand{\csym}[1]{\ensuremath{\#{c}_{\#{\hbox{\scriptsize #1}}}}}
% \newcommand{\F}{\ensuremath{\mathcal{F}}}
% \newcommand{\fsym}[2]{\ensuremath{\#{f}^{\#{\hbox{\scriptsize #1}}}_{\#{\hbox{\scriptsize #2}}}}}
% \newcommand{\rsym}[2]{\ensuremath{\#{R}^{\#{\hbox{\scriptsize #1}}}_{\#{\hbox{\scriptsize #2}}}}}
% \newcommand{\xsym}[1]{\ensuremath{\#{x}_{\#{\hbox{\scriptsize #1}}}}}
% \newcommand{\I}{\mathcal{I}}
% ********************************************************************

\usepackage{../TutTexbib/thwregex}
\usepackage{environ}
\usepackage{bm}
\usepackage{calc}
\usepackage{varwidth}
\usepackage{wasysym}
\usepackage{mathtools}

%% Tabellen
\usepackage{array}
\usepackage{multicol}

%% Bibliotheken für viele mathematische Symbole
\usepackage{amsmath, amsfonts, amssymb}



% This is a configuration file with personal tutor information.
% It is therefore excluded from the git repository, so changes in this file will not conflict in git commits.

% Copy this template, rename to config.tex and add your information below.

\newcommand{\myname}{Lukas Morawietz}
\newcommand{\mymail}{lukas.morawietz@gmail.com} % Consider using your named student mail address to keep your u**** account private.
\newcommand{\mytutnumber}{31}

% Don't forget to update ILIAS url. WARNING: Underscores '_' and Ampersands '&' have to be escaped with backslashes '\'. Blame TeX, not me.
\newcommand{\myILIASurl}{https://ilias.studium.kit.edu/ilias.php?ref\_id=855240\&cmdClass=ilrepositorygui\&cmdNode=5r\&baseClass=ilrepositorygui}

% Uncommenting this will print Socrative info with here defined roomname whenever \Socrative is called.
% (Otherwise, \Socrative will remain silent.)
% \newcommand{\mysocrativeroom}{???}

%\def\ThassesTut{}
\def\DanielsTut{}

\newcommand{\aboutMeFrame}{
	\begin{frame}{Über mich}
		\myname \\
		Informatik, 9. Fachsemester (Bachelor)
		% Lebensgeschichte...
		% Stammbaum...
		% Aufarbeitung der eigenen Todesser-Vergangenheit...
	\end{frame}
}

\def\thisyear{2019}

% Update date of exam
\def\myKlausurtermin{18.~März~2020, 14:00–16:00~Uhr}

\def\mydate#1{
		  \ifnum#1=1\relax	  23. Oktober \thisyear \
	\else \ifnum#1=2\relax	  30. Oktober \thisyear \
	\else \ifnum#1=3\relax    06. November \thisyear \
	\else \ifnum#1=4\relax    13. November \thisyear \
	\else \ifnum#1=5\relax    20. November \thisyear \
	\else \ifnum#1=6\relax    27. November \thisyear \
	\else \ifnum#1=7\relax    04. Dezember \thisyear \
	\else \ifnum#1=8\relax    11. Dezember \thisyear \
	\else \ifnum#1=9\relax    18. Dezember \thisyear \
	\else \ifnum#1=10\relax   08. Januar \nextyear \
	\else \ifnum#1=11\relax   15. Januar \nextyear \
	\else \ifnum#1=12\relax   22. Januar \nextyear \
	\else \ifnum#1=13\relax   29. Januar \nextyear \
	\else \ifnum#1=14\relax   05. Februar \nextyear \
	\else \textbf{Datum undefiniert!} 
	\fi\fi\fi\fi\fi\fi\fi\fi\fi\fi\fi\fi\fi\fi
}

\def\mylasttimestext{Was letztes Mal geschah...}

\colorlet{beamerlightred}{red!40}
\colorlet{beamerlightgreen}{green!50}
\colorlet{beamerlightyellow}{yellow!50}
\colorlet{lightred}{red!30}
\colorlet{lightgreen}{green!40}
\colorlet{lightyellow}{yellow!50}
\colorlet{fullred}{red!60}
\colorlet{fullgreen}{green}

\definecolor{myalertcolor}{rgb}{1,0.33,0.24}
\setbeamercolor{alerted text}{fg=myalertcolor}

% Flag to toggle display of KIT Logo.
% If you want to conform to the official logo guidelines, 
% you are not allowed to use the logo and should disable it
% using the following flag. Just saying.
% (But it's too beautiful, so best leave this commented. :P)
%\newcommand{\noKITLogo}{}

% Toggle handout mode by including the following line before including PraeambelTut
% and removing the % at the start (but do NOT remove the % char here, otherwise handout mode will always be on!)
% Please keep handout mode off in all commits!

% \newcommand{\handout}{}



% define custom \handout command flag if handout mode is toggled  #DirtyAsHellButWell...
\only<beamer:0>{\def\handout{}} %beamer:0 == handout mode

\newcommand{\R}{\mathbb{R}}
\newcommand{\N}{\mathbb{N}}
\newcommand{\Z}{\mathbb{Z}}
\newcommand{\Q}{\mathbb{Q}}
\newcommand{\BB}{\mathbb{B}}
\newcommand{\C}{\mathbb{C}}
\newcommand{\K}{\mathbb{K}}
\newcommand{\G}{\mathbb{G}}
\newcommand{\nullel}{\mathcal{O}}
\newcommand{\einsel}{\mathds{1}}
\newcommand{\Pot}{\mathcal{P}}
\renewcommand{\O}{\text{O}}

\def\word#1{\hbox{\textcolor{blue}{\texttt{#1}}}}
\let\literal\word
\def\mword#1{\hbox{\textcolor{blue}{$\mathtt{#1}$}}}  % math word
\def\sp{\scalebox{1}[.5]{\textvisiblespace}}
\def\wordsp{\word{\sp}}

%\newcommand{\literal}[1]{\textcolor{blue}{\texttt{#1}}}
\newcommand{\realTilde}{\textasciitilde \ }
\newcommand{\setsize}[1]{\ensuremath{\left\lvert #1 \right\rvert}}
\let\size\setsize
\newcommand{\set}[1]{\left\{#1\right\}}
\newcommand{\tuple}[1]{\left(#1\right)}
\newcommand{\normalvar}[1]{\text{$#1$}}

% Modified by DJ
\let\oldemptyset\emptyset
\let\emptyset\varnothing % proper emptyset

%\definecolor{myRed}{RGB}{255,75,20}
%\colorlet{myGreen}{KITpalegreen}

%\newcounter{tfqtempcount}
%\newcommand{\truefalseQ}[4]{
%	\setcounter{tfqtempcount}{#1}
%	\addtocounter{tfqtempcount}{1}
%	\truefalseQuestion{#1}{\value{tfqtempcount}}{#2}{#3}{#4}
%}

%\newcommand{\truefalseQuestion}[5]{\item<#1-|handout:#1-> \color<#2-|handout:#2->{#3} #4 \qquad \visible<#2-|handout:#2->{#5}}

\newcommand{\boder}{\ensuremath{\mathbin{\textcolor{blue}{\vee}}}\xspace}
\newcommand{\bund}{\ensuremath{\mathbin{\textcolor{blue}{\wedge}}}\xspace}
\newcommand{\bimp}{\ensuremath{\mathrel{\textcolor{blue}{\to}}}\xspace}
\newcommand{\bgdw}{\ensuremath{\mathrel{\textcolor{blue}{\leftrightarrow}}}\xspace}
\newcommand{\bnot}{\ensuremath{\textcolor{blue}{\neg}}\xspace}
\newcommand{\bone}{\ensuremath{\textcolor{blue}{1}}\text{}}
\newcommand{\bzero}{\ensuremath{\textcolor{blue}{0}}\text{}}
\newcommand{\bleftBr}{\ensuremath{\textcolor{blue}{\texttt{(}}}\text{}}
\newcommand{\brightBr}{\ensuremath{\textcolor{blue}{\texttt{)}}}\text{}}

\newcommand{\plB}{\plfoo{B}}
\newcommand{\plE}{\plfoo{E}}

\newcommand{\summe}[2]{\sum\limits_{#1}^{#2}}
\newcommand{\limes}[1]{\lim\limits_{#1}}

%\newcommand{\numpp}{\advance \value{weeknum} by -2 \theweeknum \advance \value{weeknum} by 2}
%\newcommand{\nump}{\advance \value{weeknum} by -1 \theweeknum \advance \value{weeknum} by 1}

\newcommand{\mycomment}[1]{}
\newcommand{\Comment}[1]{}

%% DISCLAIMER START 
% It is INSANELY IMPORTANT NOT TO DO THIS OUTSIDE BEAMER CLASS! IN ARTCILE DOCUMENTS, THIS IS VERY LIKELY TO BUG AROUND!
\makeatletter%
\@ifclassloaded{beamer}%
{
	% TODO 
	% no time...
	% redefine section to ignore multiple \section calls with the same title
}%
{
	\errmessage{ERROR: section command redefinition outside of beamer class document! Please contact the author of this code.}
}%
\makeatother%
%% DISCLAIMER END

\newcounter{abc}
\newenvironment{alist}{
  \begin{list}{(\alph{abc})}{
      \usecounter{abc}\setlength{\leftmargin}{8mm}\setlength{\labelsep}{2mm}
    }
}{\end{list}}


\newcommand{\stdarraystretch}{1.20}
\renewcommand{\arraystretch}{\stdarraystretch}  % for proper row spacing in tables

\newcommand{\morescalingdelimiters}{   % for proper \left( \right) typography
	\delimitershortfall=-1pt  
	\delimiterfactor=1
}

\newcommand{\centered}[1]{\vspace{-\baselineskip}\begin{center}#1\end{center}\vspace{-\baselineskip}}

% for \implitem and \item[bla] stuff to look right:
\setbeamercolor*{itemize item}{fg=black}
\setbeamercolor*{itemize subitem}{fg=black}
\setbeamercolor*{itemize subsubitem}{fg=black}

\setbeamercolor*{description item}{fg=black}
\setbeamercolor*{description subitem}{fg=black}
\setbeamercolor*{description subsubitem}{fg=black}

\renewcommand{\qedsymbol}{\textcolor{black}{\openbox}}

\renewcommand{\mod}{\mathop{\textbf{mod}}}
\renewcommand{\div}{\mathop{\textbf{div}}}

\newcommand{\ceil}[1]{\left\lceil#1\right\rceil}
\newcommand{\floor}[1]{\left\lfloor#1\right\rfloor}
\newcommand{\abs}[1]{\left\lvert #1 \right\rvert}
\newcommand{\Matrix}[1]{\begin{pmatrix} #1 \end{pmatrix}}
\newcommand{\braced}[1]{\left\lbrace #1 \right\rbrace}

\def\fract#1/#2 {\frac{#1}{#2}} % ! Trailing space is crucial!
\def\dfract#1/#2 {\dfrac{#1}{#2}} % ! Trailing space is crucial!

\newcommand{\Mid}{\;\middle|\;}

\let\after\circ

\def\·{\cdot}
\def\*{\cdot}
\def\?>{\ensuremath{\rightsquigarrow}}  % Fuck you, Latex

\newcommand{\tight}[1]{{\renewcommand{\arraystretch}{0.76} #1}}
\newcommand{\stackedtight}[1]{{\renewcommand{\arraystretch}{0.76} \begin{matrix} #1 \end{matrix}} }
\newcommand{\stacked}[1]{\begin{matrix} #1 \end{matrix} }
\newcommand{\casesl}[1]{\delimitershortfall=0pt  \left\lbrace\hspace{-.3\baselineskip}\begin{array}{ll} #1 \end{array}\right.}
\newcommand{\casesr}[1]{\delimitershortfall=0pt  \left.\begin{array}{ll} #1 \end{array}\hspace{-.3\baselineskip}\right\rbrace}
\newcommand{\caseslr}[1]{\delimitershortfall=0pt  \left\lbrace\hspace{-.3\baselineskip}\begin{array}{ll} #1 \end{array}\hspace{-.3\baselineskip}\right\rbrace}

\def\q#1uad{\ifnum#1=0\relax\else\quad\q{\the\numexpr#1-1\relax}uad\fi}
% e.g. \q1uad = \quad, \q2uad = \qquad etc.

\newcommand{\qqquad}{\q3uad}

\newcommand{\impl}{\ifmmode\ensuremath{\mskip\thinmuskip\Rightarrow\mskip\thinmuskip}\else$\Rightarrow$\fi\xspace}
\newcommand{\Impl}{\ifmmode\implies\else$\Longrightarrow$\fi\xspace}

\newcommand{\derives}{\Rightarrow}

\newcommand{\gdw}{\ifmmode\mskip\thickmuskip\Leftrightarrow\mskip\thickmuskip\else$\Leftrightarrow$\fi\xspace}
\newcommand{\Gdw}{\ifmmode\iff\else$\Longleftrightarrow$\fi\xspace}

\newcommand{\symbitemnegoffset}{\hspace{-.5\baselineskip}}
\newcommand{\implitem}{\item[\impl\symbitemnegoffset]}
\newcommand{\Implitem}{\item[\Impl\symbitemnegoffset]}


\newcommand{\forcenewline}{\mbox{}\\}


% proper math typography
\newcommand{\functionto}{\longrightarrow}
\renewcommand{\geq}{\geqslant}
\renewcommand{\leq}{\leqslant}
\let\oldsubset\subset
\renewcommand{\subset}{\subseteq} % for all idiots out there using subset

\newenvironment{threealign}{%
	\[
	\begin{array}{r@{\ }c@{\ }l}
}{%
	\end{array}	
	\]
}

\newcommand{\concludes}{ \\ \hline  }
\newcommand{\deduction}[1]{
	\begin{varwidth}{.8\linewidth}
		\begin{tabular}{>{$}c<{$}}
			#1
		\end{tabular}
	\end{varwidth}	
}

\definecolor{hoareorange}{rgb}{1,.85,.6}
\newcommand{\hoareassert}[1]{\setlength{\fboxsep}{1pt}\setlength{\fboxrule}{-1.4pt}\fcolorbox{white}{hoareorange}{\ensuremath{\{\;#1\;\}}}\setlength\fboxrule{\defaultfboxrule}\setlength{\fboxsep}{3pt}}

\newcommand{\mailto}[1]{\href{mailto:#1}{{\textcolor{blue}{\underline{#1}}}}}
\newcommand{\urlnamed}[2]{\href{#2}{\textcolor{blue}{\underline{#1}}}}
\renewcommand{\url}[1]{\urlnamed{#1}{#1}}

\newcommand{\hanging}{\hangindent=0.7cm}
\newcommand{\indented}{\hanging}

%requires \thisyear to be defined (s. config.tex)!
\edef\nextyear{\the\numexpr\thisyear+1\relax}


% --- \frameheight constant ---
\newlength\fullframeheight
\newlength\framewithtitleheight
\setlength\fullframeheight{.92\textheight}
\setlength\framewithtitleheight{.86\textheight}

\newlength\frameheight
\setlength\frameheight{\fullframeheight}

\let\frametitleentry\relax
\let\oldframetitle\frametitle
\def\newframetitle#1{\global\def\frametitleentry{#1}\if\relax\frametitleentry\relax\else\setlength\frameheight{\framewithtitleheight}\fi\oldframetitle{#1}}
\let\frametitle\newframetitle

\def\newframetitleoff{\let\frametitle\oldframetitle}
\def\newframetitleon{\let\frametitle\newframetitle}
% --- \frameheight constant end ---



\newenvironment{headframe}{\Huge THIS IS AN ERROR. PLEASE CONTACT THE ADMIN OF THIS TEX CODE. (headframe env def failed)}{}
\RenewEnviron{headframe}[1][]{
	\begin{frame}\frametitle{\ }
		\centering
		\Huge\textbf{\textsc{\BODY} \\
		}
		\Large {#1}
		\frametitle{\ }
	\end{frame}
}


\makeatletter
% Provides color if undefined.
\newcommand{\colorprovide}[2]{%
	\@ifundefinedcolor{#1}{\colorlet{#1}{#2}}{}}
\makeatother


\colorprovide{lightred}{red!30}
\colorprovide{lightgreen}{green!40}
\colorprovide{lightyellow}{yellow!50}
\colorprovide{beamerlightred}{lightred}
\colorprovide{beamerlightgreen}{lightgreen}
\colorprovide{beamerlightyellow}{lightyellow}
\colorprovide{fullred}{red!60}
\colorprovide{fullgreen}{green}
\definecolor{darkred}{RGB}{115,48,38}
\definecolor{darkgreen}{RGB}{48,115,38}
\definecolor{darkyellow}{RGB}{100,100,0}

\only<handout:0>{\colorlet{adaptinglightred}{beamerlightred}}
\only<handout:0>{\colorlet{adaptinglightgreen}{beamerlightgreen}}
\only<handout:0>{\colorlet{adaptinglightred}{beamerlightred}}
\only<beamer:0>{\colorlet{adaptinglightred}{lightred}}
\only<beamer:0>{\colorlet{adaptinglightgreen}{lightgreen}}
\only<beamer:0>{\colorlet{adaptinglightred}{lightred}}
\only<handout:0>{\colorlet{adaptingred}{lightred}}
\only<beamer:0>{\colorlet{adaptingred}{fullred}}
\only<handout:0>{\colorlet{adaptinggreen}{lightgreen}}
\only<beamer:0>{\colorlet{adaptinggreen}{fullgreen}}



\newcommand{\TrueQuestion}[1]{
	\TrueQuestionE{#1}{}
}

\newcommand{\YesQuestion}[1]{
	\YesQuestionE{#1}{}
}

\newcommand{\FalseQuestion}[1]{
	\FalseQuestionE{#1}{}
}

\newcommand{\NoQuestion}[1]{
	\NoQuestionE{#1}{}
}

\newcommand{\DependsQuestion}[1]{
	\DependsQuestionE{#1}{}
}

\newcommand{\QuestionVspace}{\vspace{4pt}}
\newcommand{\QuestionParbox}[1]{\begin{varwidth}{.85\linewidth}#1\end{varwidth}}
\newcommand{\ExplanationParbox}[1]{\begin{varwidth}{.97\linewidth}#1\end{varwidth}}
\colorlet{questionlightgray}{gray!23}
\let\defaultfboxrule\fboxrule

% #1: bg color
% #2: fg color short answer
% #3: short answer text
% #4: question
% #5: explanation
\newcommand{\GenericQuestion}[5]{
	\setlength\fboxrule{2pt}
	\only<+|handout:0>{\hspace{-2pt}\fcolorbox{white}{questionlightgray}{\QuestionParbox{#4} \quad\textbf{?}}}
	\visible<+->{\hspace{-2pt}\fcolorbox{white}{#1}{\QuestionParbox{#4} \quad\textbf{\textcolor{#2}{#3}}} \ExplanationParbox{#5}} \\
	\setlength\fboxrule{\defaultfboxrule}
}

% #1: Q text
% #2: Explanation
\newcommand{\TrueQuestionE}[2]{
	\GenericQuestion{adaptinglightgreen}{darkgreen}{Wahr.}{#1}{#2}
}

% #1: Q text
% #2: Explanation
\newcommand{\YesQuestionE}[2]{
	\GenericQuestion{adaptinglightgreen}{darkgreen}{Ja.}{#1}{#2}
}

% #1: Q text
% #2: Explanation
\newcommand{\FalseQuestionE}[2]{
	\GenericQuestion{adaptinglightred}{darkred}{Falsch.}{#1}{#2}
}

% #1: Q text
% #2: Explanation
\newcommand{\NoQuestionE}[2]{
	\GenericQuestion{adaptinglightred}{darkred}{Nein.}{#1}{#2}
}

% #1: Q text
% #2: Explanation
\newcommand{\DependsQuestionE}[2]{
	\GenericQuestion{adaptinglightyellow}{darkyellow}{Je nachdem!}{#1}{#2}
}

\ifnum\thisyear=2017 \else \errmessage{Old ILIAS link inside preamble. Please update.} \fi

\newcommand{\ILIAS}{\urlnamed{ILIAS}{https://ilias.studium.kit.edu/ilias.php?ref\_id=729057\&cmdClass=ilrepositorygui\&cmdNode=75\&baseClass=ilrepositorygui}\xspace}

\newcommand{\Socrative}{\only<handout:0>{socrative.com $\qquad \?> $ Student login \\ Raumname:  \mysocrativeroom\\ \medskip}}

\newcommand{\thasse}[1]{
	\ifdefined\ThassesTut #1\xspace \else\fi
}
\newcommand{\daniel}[1]{
	\ifdefined\DanielsTut #1\xspace \else\fi
}
\newcommand{\thassedaniel}[2]{\ifdefined\ThassesTut #1\else\ifdefined\DanielsTut #2\fi\fi\xspace}

\ifdefined\ThassesTut \ifdefined\DanielsTut \errmessage{ERROR: Both ThassesTut and DanielsTut flags are set. This is most likely an error. Please check your config.tex file.} \else \fi \else \ifdefined\DanielsTut \else \errmessage{ERROR: Neither ThassesTut  nor DanielsTut flags are set. This is most likely an error. Please check your config.tex file.} \fi\fi

%\newcommand{\sgn}{\text{sgn}}


% Das ist der KIT-Stil
%\usepackage{../TutTexbib/beamerthemekit}
\usepackage[deutsch,titlepage0]{../TutTexbib/KIT/beamerthemeKITmod}
\TitleImage[width=\titleimagewd]{../figures/titlepage.jpg}
%\usetheme[deutsch,titlepage0]{KIT}

% Include PDFs
\usepackage{pdfpages}

% Libertine font (Original GBI font)
\usepackage{libertine}
%\renewcommand*\familydefault{\sfdefault}  %% Only if the base font of the document is to be sans serif

% Nicer math symbols
\usepackage{eulervm}
%\usepackage{mathpazo}
\renewcommand\ttdefault{cmtt} % Computer Modern typewriter font, see lecture slides.

%% Deutsche Silbentrennung und Beschriftungen
\usepackage[ngerman]{babel}



\usepackage{csquotes}



%% Anzeigetiefe für Inhaltsverzeichnis: 1 Stufe
\setcounter{tocdepth}{1}

%% Schönere Schriften
\usepackage[TS1,T1]{fontenc}

%% Bibliothek für Graphiken
\usepackage{graphicx}

%% der wird sowieso in jeder Datei gesetzt
%%\graphicspath{{../figures/}}

%% Hyperlinks
\usepackage{hyperref}
% I don't know why, but this works and only includes sections and NOT subsections in the pdf-bookmarks.
\hypersetup{bookmarksdepth=subsection} 

%\usepackage{lmodern}
\usepackage{colortbl}
\usepackage[absolute,overlay]{textpos}
\usepackage{listings}
\usepackage{forloop}
%\usepackage{algorithmic} % PseudoCode package 

\usepackage{tikz}
\usetikzlibrary{matrix}
\usetikzlibrary{arrows.meta}
\usetikzlibrary{automata}
\usetikzlibrary{tikzmark}

% Needed for gbi-macros
\usepackage{xspace}

%%%%%%%%%%%% INHALT %%%%%%%%%%%%%%%%

%% Wochennummer
\newcounter{weeknum}



%% Titelinformationen
\title[GBI-Tutorium \mytutnumber, Woche \theweeknum]{Grundbegriffe der Informatik \\ Tutorium \mytutnumber}

\subtitle{Woche \theweeknum \ | \mydate{\theweeknum} \\ \myname \ \  \normalfont (\mailto{\mymail})}
\author[\myname]{\myname}
\institute{KIT -- Karlsruher Institut für Technologie}
\date{\mydate{\theweeknum}\ }

% Modified, DJ (better safe than sorry)
\AuthorTitleSep{ – }

%% Titel einfügen
\newcommand{\titleframe}{\frame{\titlepage}}

%% Alles starten mit \starttut{X}
\newcommand{\starttut}[1]{\setcounter{weeknum}{#1}\titleframe\frame{\frametitle{Inhalt}\tableofcontents} \AtBeginSection[]{%
\begin{frame}
	\tableofcontents[currentsection]
\end{frame}\addtocounter{framenumber}{-1}}}


\newcommand{\framePrevEpisode}{
	\begin{headframe}
		\mylasttimestext
	\end{headframe}
}

%% Legacy: Lastframe. Not for further usage!
\newcommand{\lastframetitled}[5]{
	\frame[plain]{
		\vspace{-#2pt}
		\begin{figure}[H]
			\centering
			\LARGE \textbf{\textsc{#5}} \\
			\vspace{.2\baselineskip}
			\includegraphics[scale=#1]{#3}
			\vspace{-7pt}
			\caption{ \texttt{\url{#4}} }
		\end{figure} 
	}
}
\newcommand{\lastframe}[4]{\xkcdframe{#1}{#2}{#3}{#4}{}}
\newcommand{\xkcdframe}[5]{
	\frame[plain]{
		\vspace{-#2pt}
		\begin{figure}[H]
			\centering
			\includegraphics[scale=#1]{#3}
			\vspace{-7pt}
			\caption{ \texttt{\url{#4}} }
			\vspace{5pt}
			#5
		\end{figure} 
	}
}

%% Wörter
\newcommand{\code}[1]{$\mathbf{#1}$}

%% Sterne

\newcounter{starsc}
\newcommand{\stars}[1]{
	\hfill
	\begin{minipage}{100px}
		\forloop{starsc}{0}{\value{starsc} < #1}%
		{%
			\includegraphics[scale=0.05]{star-full.pdf} \hspace*{1px}
		}%
		\forloop{starsc}{\value{starsc}}{\value{starsc} < 5}%
		{%
			\includegraphics[scale=0.05]{star-empty.pdf} \hspace*{1px}
		}
		\vspace*{2px}
	\end{minipage}
}

% No stars for me...
\renewcommand{\stars}[1]{}

\newcommand{\slideThanks}{
	\begin{frame}
		\frametitle{Credits}
		\begin{block}{}
			Vorgänger dieses Foliensatzes wurden erstellt von:\\[1em]
			Thassilo Helmold \\
			Philipp Basler \\
			Nils Braun \\
			Dominik Doerner \\
			Ou Yue \\
		\end{block}
	\end{frame}
}

%% Verbatim
\usepackage{moreverb}
\graphicspath{{../figures/}}

\begin{document}
\starttut{1}

\section{\thassedaniel{Organisatorisches}{Orga-Kram}}

\aboutMeFrame

\begin{frame}{Grundbegriffe?}
	\begin{itemize}
		\item Cloud
		\item Künstliche Intelligenz
		\item Compilerbau
		\item Big Data
		\item Apps
		\item Roboter
		\item Industrie~4.0
	\end{itemize}

	\pause
	Kommt alles noch... Aber nicht in GBI.
\end{frame}

\begin{frame}{Grundbegriffe!}
	\begin{itemize}
		\item Mengen, Relationen, Funktionen
		\item Wörter, Sprachen, Grammatiken
		\item Graphen
		\item Algorithmen
		\item Turingmaschinen
		\item Logik
		\item Prozessoren
		\item ...
	\end{itemize}
	\impl You have been warned. :P
\end{frame}

\begin{frame}[t]{Was machen wir hier?}
	\begin{itemize}[<+->]
		\item Tut-Folien bekommt ihr im \ILIAS. Keine Panik.
		\item Inhalte der VL verstehen und anwenden
		\item Beispiele, Aufgaben (da kommt \textbf{ihr} ins Spiel! \smiley)
		\item \textbf{Eure Fragen} klären! 
		\item \textbf{Kein} Ersatz für die VL. \\
			  Was zeitlich hier nicht mehr reinpasst, kann trotzdem wichtig sein!
	    \item Den Stoff etwas weniger formal behandeln \impl \thassedaniel{Formelles}{Formal-Kram} kann trotzdem wichtig sein! 
	    \item Wenn ich / meine Folien Blödsinn reden: \\
		      Schreien! / Mail an mich! \impl Sonst kann ich's nicht fixen... \frownie \\
		      Verbindlich \textbf{nur} Inhalt aus der VL/Übung! 
	\end{itemize}
	\NoQuestionE{Dürft ihr mich siezen?}{Ihr werdet sonst erschossen. \smiley \hspace{.5\linewidth}\mbox{} \\ \impl Sagt einfach \textbf{Du}.} 
\end{frame}

\begin{frame}{Leistungen}
	\begin{itemize}
		\item \textbf{Übungsschein}: \\
			  bestanden, wenn 50~\% der Gesamtpunkte aller Übungsblätter erreicht \\
		\item \textbf{Klausur}: 08.~März~2018, 14:00–16:00~Uhr \\
			  schriftlich, keine Hilfsmittel. \quad Stay tuned...
	\end{itemize}
	Wir sind Teil der \textbf{Orientierungsprüfungen} für Informatiker!
	\begin{itemize}
		\implitem Übungsschein im ersten Semester \textbf{versuchen} (da im Zweiten nicht angeboten), im Dritten bestehen! 
		\implitem Klausur: \\
		Teilnahme auch ohne Übungsschein möglich. \\
		Versuch \textbf{spätestens} im zweiten Semester, bestehen im Dritten!
	\end{itemize}
	
	\pause
	\DependsQuestionE{Ich studiere \enquote{Europäische Kultur und Ideengeschichte},\\ brauche ich den Übungsschein?}{Es kann nie schaden, die Blätter zu machen. Außerdem: Wer die Blätter macht, fällt seltener durch. Wirklich! \smiley}
\end{frame}

\begin{frame}{Übungsblätter}
	\begin{itemize}
		\item \textbf{Ausgabe}: Mittwochs, alle zwei Wochen im \ILIAS \\
			  \textbf{Abgabe}: Donnerstag zwei Wochen später, 16:00~Uhr, im GBI-Kasten im U-Geschoss des Infobaus
		\item Handschriftlich, \\ 
		mit \textbf{Deckblatt} (Tutoriums-Nr. \mytutnumber\  drauf!), \\ 
		\textbf{getackert} abgeben!
		\item Rechtzeitig! \\(16:01~Uhr ist höchstwahrscheinlich zu spät!)
		\item \textbf{Keine} Plagiate, \textbf{keine} Gruppenabgabe (sonst \textbf{Übungsschein weg}!) 
		\item Voraussichtlich 6--7~Blätter dieses Semester
		\item \textbf{Rückgabe}: Hier im Tut \\ (Nach zwei Wochen Abholung beim Stüker/Zenkel... tba. Keine Ahnung.) %TODO
	\end{itemize}
	
\end{frame}

\begin{frame}{Übungsblätter}
	Außerdem:
	\begin{itemize}
		\item Text in Schwarz oder Blau (oder andere dunkle Farbe $\notin \{\text{Rot}, \text{Grün}, \text{Weiß}\}$ ) \\
			  \textbf{Kein Bleistift!}
		\item Links und rechts bitte einen Rand frei lassen. \\
			  Dann kann ich's auch korrigieren. \smiley
		\item \textbf{Welche Aufgabe} ist das gerade?! \impl Deutlich machen!
		\item Was ich nicht lesen kann, gibt keine Punkte! \\
			  \small Was ich schwer lesen kann, u.~U. auch... 
	\end{itemize}
\end{frame}



\begin{frame}{Fragen? Hilfe?}	
	\begin{block}{Fragen}
		\begin{itemize}
			\item \textbf{Hier} im Tut!
			\item Ins \ILIAS. (Dann haben alle was davon. \smiley) \\
				  Im ILIAS hat's sogar einen Sprechstunden-Chat! :D
			\item Organisatorischer Spezialkram? \impl an Zenkel/Stüker (\mailto{thomas.zenkel@kit.edu} / \mailto{sebastian.stueker@kit.edu}) \\
			Bei Orga-Problemen aber \textbf{Name} und \textbf{Matrikelnummer}  mitangeben! \\
			Fachliches bitte nicht per Mail, lieber \impl Sprechstunde, ILIAS...!
			\item Tut-spezifisches an \textbf{mich} (\mailto{\mymail})
		\end{itemize}
	\end{block}
	
	\begin{block}{Content}
		\begin{itemize}
			\item VL-Folien, Skript, Ü-Blätter, Altklausuren...: im \ILIAS.
			\item Archiv: \url{http://gbi.ira.uka.de} \\
			Weitere Altklausuren, Alte ÜBs, Lösungen...
		\end{itemize}
	\end{block}
\end{frame}


\section{Information}

\begin{frame}{Nachricht, Information, ...}
	
	\begin{block}{Nachricht}
		      Mitteilung, bei der vom Medium und den Einzelheiten der Signale abstrahiert wird.
	\end{block}
	
	\pause
	\begin{block}{Information}
		      Bedeutung, die einer Nachricht zugeordnet wird (kontextabhängig!).
	\end{block}
	
	\pause
	\begin{block}{Informationsgehalt}
		Bei gleicher Wahrscheinlichkeit: Anzahl der Elemente... \\
		      Naturalis: log$_{e}$ \\
		      Hartley: log$_{10}$ \\
		      Shannon: log$_{2}$ \\[1em]
		\impl Mehr dazu in TGI
	\end{block}

\end{frame}

\section{Mengen}

\morescalingdelimiters

\begin{frame}{Mengen}
	\pause % -> HIER Sammeln, was Tutanden bereits wissen (an der Tafel)
	\begin{block}{Definition: Mengen} \vspace{-.5\baselineskip}
		\begin{itemize}
			\item \textbf{Menge} $M$: Ansammlung verschiedener Objekte
			\item $m \in M$: Objekt $m$ ist \textbf{Element} der Menge $M$
			\item Schreibweisen: \\
			$M = \{m_1, m_2, m_3 \}, \qquad M = \{m \mid m\text{ ist toll} \}$
			\item Reihenfolge der Aufzählung wurscht, Duplikate auch: \\
			$\{2,1,3,1,4\} = \{1,2,3,4\}$
			\item \textbf{Leere Menge} $\emptyset$: enthält keine Elemente, \quad also $\{\} = \emptyset$
		\end{itemize}
	\end{block}
	\pause
	
	\begin{exampleblock}{Beispiel}
		Wichtige Mengen sind \\ 
		\centered{$\Z, \Q, \R, \C, \qquad \N_+ = \{1, 2, 3, \dots\}, \quad \N_0 = \{0, 1, 2, 3, \dots\}$} 
		Es gilt:\\ 
		\centered{$-5 \in \Z\qquad {-5} \notin \N_0$}
	\end{exampleblock}
	
\end{frame}

\mycomment{ % Too complicated and too little time.
	\begin{frame}{Anwendung}
		Wir haben ein „Universum“ an Dingen:\\ 
		$U = \{ $ Gurken, Obst, Fußball, Hockey, Äpfel, Maracujas, Tomaten, Bananen, Karotten, Wurzelpetersilie, Sport,  Orangen, Gemüse, Basketball, Ersti-Weitwurf $ \} $  \\[0.5em]
		
		\pause
		Ordnen wir diese in eine Teilmenge $C$ für Kategorien und jeweils eine Teilmenge $E_c$ für die Elemente einer Kategorie $c \in C$. \\ {\small (\textbf{Überhaupt} nicht konstruiert!)} \\[0.5em]
		
		\pause
		$C = \{$ Obst, Gemüse, Sport $\}$ \\[0.3em]
		$E_{\text{Obst}} = \{$ Äpfel, Orangen, Maracujas, Bananen $\}$ \\
		$E_{\text{Gemüse}} = \{$ Tomaten, Karotten, Gurken, Wurzelpetersilie $\}$ \\
		$E_{\text{Sport}} = \{$ Fußball, Basketball, Hockey, Ersti-Weitwurf $\}$ \\
		
		\vspace{.5\baselineskip}
		\Large \alert{
			\textbf{Achtung}: 
			$$\underbrace{\text{Obst}}_{\mathllap{\text{ein Bezeichner, keine Menge}}}
			\neq 
			\underbrace{E_{\text{Obst}}}_{\mathrlap{\text{eine Menge}}}$$
		}
		
	\end{frame}
}


\begin{frame}{Kardinalität, Teilmengen, Gleichheit}
	
	\begin{block}{Definitionen}  \vspace{-.4\baselineskip}
		\begin{itemize}
			\item \textbf{Kardinalität} $\setsize{M}$: \\
			Anzahl der Elemente einer endlichen Menge $M$ \\
			\pause
			\item $N \subseteq M$: \\
			Menge $N$ ist \textbf{Teilmenge} von $M$, also jedes Element aus $N$ auch in $M$ \\
			Es gilt: \qqquad $ N \subseteq M \iff \forall n \in N : n \in M$
			\pause
			\item $N = M$: \\
			Mengen $N$ und $M$ sind \textbf{gleich}, enthalten also die gleichen Elemente \\
			Es gilt: \qqquad $ N = M \iff N \subseteq M \ \text{und} \ M \subseteq N$
		\end{itemize}
	\end{block} 
\end{frame}

\begin{frame}{Kardinalität, Teilmengen, Gleichheit}
	
	\begin{exampleblock}{Beispiel}
		
		\begin{align*}
		\setsize{\{1,2,3,2,1\}} &= \only<2->{3} \\
		\setsize{ \emptyset } &= \only<3->{0} \\
		\setsize{\{\{\}, 1, \{2,3\} \}} &= \only<4->{3} \\
		\{1,2\} \only<1-4|handout:0>{ &\mathrel{?\ }  } \only<5->{&\subseteq}  \{1,2,3\} \\
		\{1,2\} \only<1-5|handout:0>{ &\mathrel{?\ } } \only<6->{&\nsubseteq} \{\text{Hund}, \text{Katze}, \text{Maus}\} \\
		\{1,2,3\} \only<1-6|handout:0>{ &\mathrel{?\ } } \only<7->{&\supseteq} \{1,2\}
		\end{align*}
	\end{exampleblock} 
	

\end{frame}

\begin{frame}{Schnitt und Vereinigung}
	
	\begin{block}{Definition}
		Seien $M$ und $N$ zwei Mengen, so heißen
		$$M \cap N := \left\{x \Mid x \in M \text{ und } x \in N \right\} \qquad M \cup N := \left\{x \Mid x \in M \text{ oder } x \in N \right\} $$
		\textbf{Durchschnitt} bzw. \textbf{Vereinigung} von $M$ und $N$.\\[1em] 
		\pause
		$M$ und $N$ heißen \textbf{disjunkt}, wenn sie keine gemeinsamen Elemente haben, also $M \cap N = \emptyset$ (der Schnitt ist leer).
	\end{block}
	
	\pause
	\begin{exampleblock}{Beispiel}
		\centered{$ \{1,2\} \cup \{2,3\} = \{1,2,3\} \qquad \{1,2\} \cap \{2,3\} = \{2\} $}
	\end{exampleblock}
	
	\pause
	
	\begin{block}{Lemma}
		Es gilt: \\ \centered{$\setsize{M \cup N} \ = \ \setsize{M} + \setsize{N} - \setsize{M \cap N}$}
	\end{block}
	
	
\end{frame}

\begin{frame}{Mengendifferenz}
	\begin{block}{Definition}
		Seien $A$ und $B$ zwei Mengen, so heißt $$ A\setminus B := \left\{ x \Mid x\in A \text{ und } x\notin  B  \right\} $$ die \textbf{Differenz}(-menge) von $A$ und $B$.
	\end{block}
	
	\pause
	
	\begin{block}{Lemma}
		Es gilt: $$ A \subseteq B \iff A \setminus B = \emptyset $$
	\end{block} 
	
	\pause
	
	\begin{exampleblock}{Weitere Beispiele}
		\vspace{-2\baselineskip} % WtF LaTeX??
		\begin{align*}
		A \cup \emptyset &= \uncover<4->{  A }  \\
		A \cap \emptyset &= \uncover<5->{ \emptyset }\\
		\N_+ \cup \{0\} &= \uncover<6->{ \N_0} \\
		\N_0 \setminus \{0\} &= \uncover<7->{ \N_+} 
		\end{align*}
	\end{exampleblock}
	
\end{frame}

\begin{frame}[t]{Mengengleichheit: Beispiel}
	\begin{itemize}
		\item<1-> Sei $ A $ und $M$ beliebige Mengen. Zeigt, dass gilt 
		\begin{align*}
		A &=  \underbrace{ \left(A \setminus M \right)}_{T_1} \cup  \underbrace{ \left(A \cap M \right)}_{T_2}  
		\end{align*}		 
		\only<2-5>{
			\item<2-5> \textbf{Richtung}: $ A \subseteq T_1 \cup T_2 $ 
			\item<3-5> Wähle $x\in A$ und unterscheide:
			\item<4-5> Fall 1: Ist $x\in M$, so gilt $x\in A $ und $x\in M$ und damit $x\in A\cap M  = T_2 $
			\item<5-5> Fall 2 : Ist $x\notin M$, so gilt $x\in T_1$, da $T_1 = \left\{x \Mid x \in A \text{ und } x\notin M \right\} $  }
		\only<6->{\item<6-> \textbf{Richtung}: $T_1 \cup T_2 \subseteq A $ 
			\item<7-> Wähle $x\in T_1 \cup T_2$. Dies bedeutet $x\in T_1 $ oder $x\in T_2$. 
			\item<8-> Fall 1: $x\in T_1$ (angenommen $T_1 \neq \emptyset$). Aus Definition folgt $x\in A$.
			\item<9-> Fall 2: $x\in T_2$ (angenommen $T_2 \neq \emptyset$). Somit $x\in A$ und $x\in M$. $\qed$}
	\end{itemize}
\end{frame}

\begin{frame}{Eine Menge Mengen...}
	\begin{block}{Aufgabe}
		Wir haben Mengen $A = \{1, 2\}, B = \{3\}, C = \{1, 3\}; A,B,C \subseteq M = \{1, 2, 3\}$.\\
		Dann ist:
		\begin{align*}
		A \cup B &= \visible<2->{ \{1, 2, 3\} }  \\
		A \cap C &= \only<3->{ \{1\} }\\
		A \setminus C &= \only<4->{ \{2\} }\\
		B \setminus A &= \only<5->{ \{3\} }\\
		A \cup (B \setminus C) &= \only<6->{ \{1, 2\} }\\
		C &= \{1, 3\} \\
		(A \setminus C) \cup B &= \only<7->{ \{2, 3\} }\\
		A \cap B &= \only<8->{ \emptyset }
		\end{align*}
	\end{block}
\end{frame}



\section{Potenzmengen}

\begin{frame}{... in einer Menge! (Potenzmengen)}
	\begin{block}{Definition}
		Die \textbf{Potenzmenge} $2^M$ oder auch $\Pot (M)$ ist die Menge aller möglicher Teilmengen von $M$. 
		\begin{align*}
			2^M := \left\{ N \Mid  N \subseteq M \right\}
		\end{align*}
	\end{block}
	\pause
	
	\begin{exampleblock}{Beispiel}
		Sei $M = \left\{ 1,2,0 \right\} $. \\ \pause
		
		Dann gilt  
		\begin{align*}   
		2^M &= \left\{ \emptyset, \{ 0 \}, \{ 1 \}, \{ 2 \}, \{ 0,1 \} , \{ 0,2 \}, \{ 1,2 \}, \{ 0,1,2 \} \right\}
		\end{align*}
		
		\textbf{Beachte}: Es gilt immer $M \in 2^M \ \text{und} \ \emptyset \in 2^M$.
	\end{exampleblock}
	
\end{frame}

\begin{frame}{Potenzmengen}
	\begin{block}{Aufgabe}
		Wie viele Elemente enthält $2^M$? \\[0.5em]
		
		\pause
		$2^{\setsize{M}}$
	\end{block}
	
	\pause
	\begin{block}{Noch ne Aufgabe}
		Gebt eine Abbildung $\phi \colon 2^{M} \functionto 2^{M}$ so an,
		dass für jedes $L \in 2^{M}$ und für jedes $w \in M$ gilt:
		\begin{equation*}
			w \in L \text{ genau dann, wenn } w \notin \phi(L).
		\end{equation*}
		
		\pause
		\begin{threealign}
			\phi \colon 2^{M} &\functionto& 2^{M},\\
			L &\mapsto& M \setminus L.
		\end{threealign}
	\end{block}
\end{frame}

\section{Paare}

\begin{frame}{Paare}
	\begin{block}{Definition}
		Seien $A$ und $B$ zwei Mengen und $a \in A$, $b \in B$.\\
		$$(a, b)$$ heißt \textbf{Paar} mit der ersten Komponente $a$ und der zweiten Komponente $b$.\\[1em]
		\pause
		Paare sind \textbf{keine Mengen, sondern was \alert{völlig anderes}!} \\
		In Paaren sind Duplikate \textbf{möglich} und die Reihenfolge der Elemente ist \textbf{wichtig}.\\
	\end{block}

	\pause
	\begin{exampleblock}{Beispiel}
		$$ (42, \triangle) = (42, \triangle) \qquad (42, \triangle) \neq (\triangle, 42) \qquad (1, 1) \neq \underbrace{(1)}_{\mathclap{\text{überhaupt gar kein Paar!}}} $$
	\end{exampleblock}
\end{frame}

\begin{frame}
	\centering
	\Huge
	\alert{Runde und geschweifte Klammern 
		   $$\Big(\;\Big) \textbf{\quad vs. \quad} \Big\{\;\Big\}$$ \\
		   \textbf{NICHT VERWECHSELN!}}
\end{frame}

\begin{frame}{Paare}
	Das Konzept der Paare lässt sich auf das Konzept der Mengen zurückführen.
	
	\begin{block}{Aufgabe}
		Gegeben sei die Menge $M = \{m_1, m_2\}$.\\
		Wie kann man die Paare $(m_1, m_2)$ und $(m_2, m_1)$ eindeutig darstellen, allerdings nur unter Verwendung von Mengen und $m_1, m_2$?
	\end{block}
	\pause
	
	\begin{block}{Lösung}
		Wir definieren: \\ 
		$(m_1, m_2) := \left\lbrace m_1, m_2, \{m_1\}\right\rbrace$ und $(m_2, m_1) := \left\lbrace m_1, m_2, \{m_2\}\right\rbrace$
	\end{block}
\end{frame}

\begin{frame}
	\begin{block}{Was ihr nun wissen solltet}
		\begin{itemize}
			\item wie man mit Mengen umgeht
			\item wie man mit noch mehr Mengen umgeht
			\item was Paare sind
		\end{itemize}
	\end{block}

	\begin{block}{Was nächstes Mal kommt}
		\begin{itemize}
			\item wie man mehr als zwei Elemente geordnet zusammenbringt (Tupel)
			\item wie man mit Relationen Ordnung in das Chaos bringt
			\item[...] und noch vieles mehr!
		\end{itemize}
	\end{block}
\end{frame}

\only<handout:0>{\slideThanks}

%% Letzte Seite 
\lastframe{0.35}{30}{xkcd/good_code.png}{http://xkcd.com/844}

\only<beamer:0>{\slideThanks}

\end{document}
